\documentclass[12pt]{article}
\usepackage[utf8]{inputenc}
\usepackage[T1]{fontenc}
\usepackage{hyperref}
\title{Pin Language Resources for K-12 Curriculum Development}
\author{Compiled by Claude Sonnet 5.5 (claude-sonnet-5-5) for Andr\'as Kornai}
\date{}
\begin{document}
\maketitle

\section{Introduction}
Pin (ISO 639-3 wmx) is also known as Womo. It is a small Skou-family language of Sandaun Province, Papua New Guinea, spoken around Onei village. Estimates of the number of speakers vary widely. Wikipedia gives 6 native speakers (2016), Global Recordings Network gives 190, and Joshua Project gives figures under 1,000 and notes possible assimilation into neighbouring groups. Womo and Sumararu are sometimes treated as one language. The search found no evidence of school use.

\section{Primary Educational Resources}
No primers, textbooks or school programmes were found.

\section{Linguistic and Reference Materials}
\begin{itemize}
\item Wikipedia, Womo language: location, speaker counts, dialects and classification in the Skou family. \url{https://en.wikipedia.org/wiki/Womo_language}
\item WALS entry for Womo. \url{https://wals.info/languoid/lect/wals_code_wom}
\end{itemize}
No dictionary or grammar was found in the search results.

\section{Religious and Community Resources}
Joshua Project reports no Bible translation, New Testament or recordings for this group. \url{https://joshuaproject.net/people_groups/20773}

\section{Digital and Audio Resources}
Global Recordings Network has a language page for wmx. The search did not show what recordings, if any, it holds. \url{https://globalrecordings.net/en/language/wmx}

\section{Curriculum Development Context}
The language is severely endangered, and the very small number of speakers makes documentation the priority before any curriculum work. No pedagogical materials were found.

\end{document}